% TODO
% Response to referee 2, comment 2
Plan:
\begin{enumerate}
\item This is true.
\item The more proper statement is that the Higgs mode should modulate all of the non-background SHG susceptibility elements equally.
\item The statement on page 2 says that it should APPEAR in all four polarization channels, not that it should appear in all polarization channels EQUALLY.
\item The transverse mode ONLY affects PS/SP.
This is not consistent with a Higgs mode that modulates the non-background components equally.
\item Ajesh it would be nice to write some pretty equations to indicate what the effect of a background contribution is.
\item For the referee's second point, I think they are talking about the same thing as referee 1 where this contribution is order $(\delta |P|)^2$?
But honestly I'm really confused about this whole comment!

\item OK this is what our paper said:
"According to Eq. (1), the Higgs mode should therefore modulate all of the elements of $\chi_{ijk}$ equally, and should thus appear in all four SHG polarization channels".
This statement is of the form "If A, then B"
But we can make a weaker version of A and still conclude B
B is the important part
Thus "If weaker version of A, thus B" is still a proper thing to say in our paper and doesn't change any conclusions
We thank the referee for pointing out this confusion

\item Also eq. 1 reduces $\chi_{ijkl}$ to $\chi_{ijky}$ which is untrue as the spins tilt out of plane by some degree.
This doesn't change any conclusions though we just have to make sure we fix it.

\item OK so here's our argument for why a weaker version of A is sufficient: if the different modulations are potentially different, then how do we even know which one is the higgs and which one is the transverse? Well, if they were switched, then it would imply an exact cancelation between the background and the slow mode, but we know that can't be possible.
\textcolor{red}{Question: how much detail should we go into?}
\end{enumerate}
